\chapter{Conclusion and Future Work}
\label{ch:conclusion}
\sloppy
\hbadness=8000

\section{Research Summary}
This thesis designed, developed, and empirically evaluated a machine learning-based fabric consumption forecasting system for the apparel manufacturing sector, with the objective of improving upon the systematic inaccuracies of the traditional Bill of Materials estimation formula. The research was structured across six sequential phases: (1)~a literature review that identified the gap between the BOM formula's static parameters and the dynamic realities of production; (2)~physics-informed synthetic dataset generation producing 5{,}000 production orders spanning January~2022 to September~2024; (3)~a four-stage data preprocessing and feature engineering pipeline; (4)~training, evaluation, and ensemble construction across four machine learning architectures---Linear Regression, Random Forest, XGBoost, and LSTM; (5)~segmented performance analysis across four production seasons and five fabric categories; and (6)~deployment of all trained models within a six-module Streamlit web application designed for direct industrial use. All experiments were conducted with a fixed random seed of~42 and a strict 64/16/20 train/validation/test split, ensuring that all reported performance figures are drawn exclusively from data unseen during any stage of model training or hyperparameter selection.

\secdiv
\section{Key Findings}
The experimental results establish a consistent performance hierarchy across the four models. Linear Regression, as the interpretable baseline, performs substantially below the ML alternatives (RMSE 2{,}062\,yd, MAPE 49.28\%, R\textsuperscript{2}=0.8669), confirming that the relationship between order attributes and fabric consumption is strongly non-linear and cannot be adequately represented by a linear combination of features. Random Forest improves markedly over LR (RMSE 821\,yd, MAPE 15.92\%, R\textsuperscript{2}=0.9789), demonstrating that ensemble tree methods handle the non-linear interactions effectively, though significant residual errors remain at large order scales. XGBoost achieves the strongest classical performance (RMSE 574\,yd, MAPE 8.14\%, R\textsuperscript{2}=0.9897), with the regularised sequential boosting approach systematically focusing successive trees on the hardest-to-predict residuals. The LSTM achieves the lowest RMSE of all models (346\,m / 379\,yd), the lowest MAPE (5.14\%), and the highest R\textsuperscript{2} (0.9955)---establishing it as the best single model and demonstrating a 21.8\% RMSE improvement over the BOM physics baseline (484\,yd). The dynamic inverse-RMSE\textsuperscript{2} ensemble (RMSE 422\,yd, MAPE 5.92\%, R\textsuperscript{2}=0.9944) marginally underperforms the LSTM in point accuracy but provides enhanced robustness for procurement-critical decisions by combining models with diverse error modes.

XGBoost feature importance analysis identified garment type (importance score 0.518) and order quantity (0.297) as the two dominant predictors, jointly accounting for 81.5\% of the model's decision weight. This finding validates the BOM formula's core structure while explaining its systematic deviation: the static BOM coefficients cannot capture the non-linear interaction between garment complexity and order scale that the gradient boosting model learns from data. The variance analysis revealed a mean BOM bias of only $-$0.65\% across all orders, but with a range from $-$14.22\% to $+$19.92\% and a strong directional dependence on order size: the BOM over-consumes relative to plan for small orders (100--500 units: $+$4.49\%) and under-consumes for large orders (3{,}000--5{,}000 units: $-$2.79\%).

Segmented analysis revealed meaningful heterogeneity in forecasting difficulty. Winter orders proved the hardest to forecast (MAPE 9.67\%), exhibiting 37\% higher error than spring orders (MAPE 7.05\%), attributable to the higher frequency of multi-layer and outerwear garments in winter production. Denim fabric was the most predictable category (MAPE 6.27\%), while polyester orders exhibited the highest residual error (MAPE 10.31\%)---a 64\% differential reflecting polyester's more variable shrinkage behaviour across dyeing conditions. Five-fold cross-validation on the combined train/validation set confirmed that XGBoost's generalisation is stable, with a cross-validation RMSE standard deviation of only 28.3\,yd (4.7\% of the 598\,yd CV mean). Confidence interval coverage testing confirmed that all four models achieve nominal 90\% coverage on the test set, with LSTM providing the tightest intervals ($\pm$7.7\% half-width).

\secdiv
\section{Contributions to Knowledge}
This thesis makes six contributions that advance the literature on machine learning for apparel material planning.

\textbf{C1: First four-way algorithm comparison for fabric consumption forecasting.} Prior studies in the apparel ML literature typically evaluate one or two models. This thesis provides a controlled four-way comparison of Linear Regression, Random Forest, XGBoost, and LSTM on the same dataset under identical preprocessing and evaluation conditions, establishing for the first time the relative advantage of deep learning (LSTM) over gradient boosting and ensemble methods for this specific regression task.

\textbf{C2: Open physics-informed synthetic dataset generator.} The dataset generation module produces realistic, physics-constrained production orders from a parameterised BOM formula with calibrated noise. The generator is fully reproducible (seed~42) and extensible to other garment categories, fabric types, and factory configurations, providing the research community with a public benchmark for future comparison studies.

\textbf{C3: First Streamlit deployment with confidence intervals, batch processing, and ROI estimation.} Existing apparel forecasting research is conducted in academic experimental environments. This thesis bridges the research-to-practice gap by deploying all trained models in a six-module Streamlit application that provides single-order prediction, batch CSV processing, 90\% confidence intervals, cross-validation visualisation, and an integrated ROI calculator---all accessible via a standard web browser without specialist data science tooling.

\textbf{C4: Quantification of seasonal and fabric-type performance heterogeneity.} The 37\% seasonal MAPE differential (winter vs.\ spring) and 64\% fabric-type MAPE differential (polyester vs.\ denim) are reported here for the first time for a fabric consumption regression task, providing procurement planners with empirical guidance on where ML predictions should be treated with greater caution.

\textbf{C5: Yards-based performance benchmark.} Published apparel forecasting studies report accuracy in percentage terms or in metric units. This thesis additionally reports RMSE in yards---the operational unit used by procurement teams for fabric ordering---making the performance figures directly interpretable by industry practitioners without unit conversion.

\textbf{C6: Quantified multi-dimensional impact assessment.} The economic, operational, and environmental implications of ML adoption are reported with explicit computational basis: the \$180{,}000 annual saving, sub-four-month payback, \$730{,}000 five-year NPV, and 6{,}850\,kg CO\textsubscript{2} avoidance per 1{,}000 orders are all derived from the empirically validated 21.8\% RMSE improvement rather than from assumed savings rates.

\secdiv
\section{Limitations of the Study}
Four limitations bound the interpretation of these findings and should be explicitly acknowledged.

\textbf{Synthetic dataset.} The most significant limitation is that all models were trained and evaluated on a physics-informed synthetic dataset rather than on live factory data. While the dataset was designed to replicate the statistical properties and directional biases of real BOM estimation, the noise model and parameter distributions were calibrated from published industry benchmarks rather than from a specific factory's production records. Real factory data would introduce additional sources of variance---supplier substitution, last-minute design changes, fabric defect rates varying by batch---that are not captured in the synthetic generator. The performance improvements reported here should therefore be treated as indicative rather than audited projections until validated on live procurement data.

\textbf{Single-timestep LSTM architecture.} The LSTM was implemented as a single-timestep regressor operating on tabular feature vectors rather than as a sequence model processing the time-ordered history of a factory's production schedule. While this formulation achieves the highest test-set accuracy among the four models, it does not exploit the temporal dependencies between successive production orders that a genuine sequence LSTM could learn. A multi-step LSTM consuming the last $k$ orders' residuals as additional features may improve forecasting accuracy, particularly for the high-variance winter season.

\textbf{Absent features: GSM and shrinkage coefficient.} Fabric weight (grams per square metre, GSM) and material-specific shrinkage coefficients are absent from the nine-feature input set, as they were not available in the dataset specification. Both variables are known from textile engineering to be significant determinants of cutting consumption: heavier fabrics require adjusted marker efficiencies, and fabrics with high shrinkage rates require additional allowance. Their inclusion in a future dataset and model could substantially reduce the polyester MAPE, which exhibited the highest residual error in this study.

\textbf{Approximate economic benchmarks.} The economic impact projections use dataset-derived per-order waste-cost values (based on fabric-specific unit prices calibrated from published apparel industry references) and a one-time implementation cost of \$50{,}000 that are representative of mid-scale RMG factory contexts but are not sourced from a specific factory's audited financial records. Actual savings will vary with contracted fabric prices, factory IT infrastructure maturity, and the volume of orders processed annually. The projections are presented as a framework for impact estimation rather than as audited financial forecasts.

\secdiv
\section{Recommendations for Practice}
Based on the experimental findings, five operational recommendations are offered for apparel manufacturing factories considering ML-based fabric consumption forecasting.

\textbf{Model selection.} For the majority of production orders, the LSTM model should be the primary forecasting tool, given its superior test-set accuracy (RMSE 379\,yd, MAPE 5.14\%, R\textsuperscript{2}=0.9955). For procurement-critical orders where the cost of a fabric shortfall is high---such as single-source luxury fabric orders or time-sensitive export orders---the ensemble model is recommended instead, as its averaging mechanism reduces the probability of large individual prediction errors even at a small accuracy cost.

\textbf{Feature prioritisation.} Garment type and order quantity account for 81.5\% of XGBoost feature importance. Factories collecting limited data should ensure these two attributes are recorded with high accuracy and consistency before investing in the collection of secondary attributes such as fabric finish or complexity rating.

\textbf{Seasonal recalibration.} The 37\% MAPE differential between winter and spring suggests that model performance is meaningfully season-dependent. Factories should retrain models at the start of each production season by incorporating the previous season's actual consumption data, using the single-command retraining capability (\texttt{python train\_models.py}). This periodic recalibration prevents performance degradation as garment trends and seasonal product mixes evolve.

\textbf{Data enrichment.} Factories should prioritise the collection of fabric GSM and material shrinkage coefficients as part of their procurement data standards. The 64\% MAPE differential between polyester and denim strongly suggests that these physical properties drive a substantial portion of the residual error for synthetic fabric orders, and their inclusion in future model iterations is expected to produce the largest marginal accuracy improvement.

\textbf{Graduated adoption.} Factories new to ML-based procurement should begin by running the ML prediction in parallel with the existing BOM estimate as a second-opinion tool, comparing both forecasts with actuals over at least one production season before shifting procurement decisions to the ML recommendation. The confidence interval display in the Streamlit application facilitates this comparison by making model uncertainty transparent and preventing over-reliance on point estimates.

\secdiv
\section{Future Research Directions}
Seven directions are identified for extending this research.

\textbf{Real-world validation.} The most important next step is to deploy the system in a live factory environment and retrain the models on actual production records. Real-world validation would quantify the gap between synthetic-dataset performance and production-environment performance, and would provide the empirical foundation for a peer-reviewed impact study using audited financial and environmental data.

\textbf{Stacking ensemble.} The current ensemble uses fixed inverse-RMSE\textsuperscript{2} weights computed from validation-set performance. A learned stacking ensemble---in which a trained meta-learner (e.g., Ridge Regression or a shallow neural network) combines base-model predictions using cross-validated training signals---would adaptively weight models based on input characteristics rather than applying a single global weighting, and is expected to narrow the gap between the ensemble and the best single model on the held-out test set.

\textbf{Sequence LSTM with temporal order history.} Reformulating the LSTM as a genuine sequence model that consumes the temporal history of a factory's production orders---rather than treating each order as an independent tabular observation---would allow the model to capture trends such as within-season demand ramp-up and supplier-specific fabric utilisation drift. This formulation would require a longer historical dataset than was available for this study.

\textbf{Fabric-specific sub-models.} The 64\% MAPE differential between polyester and denim suggests that a portfolio of fabric-specific models, each trained on orders for a single fibre category, may outperform the single global model for the highest-variance categories. This approach trades data volume per model for reduced within-category heterogeneity.

\textbf{Transformer architecture.} Transformer-based sequence models have demonstrated state-of-the-art performance across a range of industrial forecasting tasks. Applying a Transformer architecture to the production order sequence, with self-attention over historical orders as the input representation, could improve winter-season forecasting accuracy where temporal patterns between successive heavy-fabric orders are most pronounced.

\textbf{Conformal prediction.} The current confidence intervals are derived from the empirical residual distribution of each model. Conformal prediction provides a statistically rigorous framework for constructing coverage-guaranteed intervals that are valid under minimal distributional assumptions, without requiring the Gaussian residual assumption implicit in the current quantile-based approach. Implementing conformal prediction wrappers around each model would strengthen the statistical validity of the reported 90\% coverage figures.

\textbf{IoT and real-time integration.} Integrating the forecasting system with real-time data from spreading machines, CAD marker planning software, and ERP fabric issue records would enable the model to incorporate actual cut-plan efficiency and real-time yield information rather than relying on planned pattern dimensions. IoT connectivity would also enable the detection of anomalous consumption events in near-real-time, triggering automatic re-procurement alerts before production stoppages occur.

\secdiv
\section{Concluding Remarks}
This thesis has demonstrated that machine learning---specifically the LSTM model and a dynamic weighted ensemble---can materially improve upon the BOM formula for fabric consumption forecasting in apparel manufacturing. The LSTM's 21.8\% RMSE reduction over the BOM baseline (379\,yd vs.\ 484\,yd) is not merely a statistical result: translated to operational and financial terms, it represents \$180{,}000 in annual waste cost savings, a payback period of under four months on a \$50{,}000 implementation investment, a five-year net present value of \$730{,}000, and 6{,}850\,kg of avoided CO\textsubscript{2} emissions per 1{,}000 production orders. The Streamlit deployment bridges the gap between academic model development and industrial practice, providing procurement teams with a browser-accessible forecasting tool that requires no specialist data science expertise to operate.

The limitations of synthetic data, absent physical fabric attributes, and single-timestep LSTM formulation define a clear research agenda whose central priority is real-world validation with live factory data. Until that validation is completed, the quantitative projections in this thesis should be treated as carefully derived estimates rather than audited outcomes. Nevertheless, the methodological framework, the open dataset generator, the validated model architectures, and the deployed application together constitute a replicable, extensible foundation upon which factory-specific forecasting systems can be built. The transition from formula-based to data-driven fabric procurement is both technically feasible and economically compelling; this thesis provides the evidence and the artefacts to make that transition.
